% =============================================================================
% 1 Introduction. Written 2026-10-20 from Ben's structure, completed and checked
% in /mnt/beegfs/benishay/paper_versions/sec0-2-outline.md and approved by Ben.
% Every claim about another system carries a % src: with its quote and section.
% Three passages of an earlier draft were copied from DirectKV (abstract, s1)
% and DAK (s1); this text replaces them and must not drift back.
% =============================================================================

\section{Introduction}\label{sec:intro}

LLM inference has two phases with opposite profiles. Prefill processes the
whole prompt at once and is compute-bound. Decode emits one token per request
per step, and every step reads all of the model's weights and every running
request's KV cache. Decode therefore does few operations per byte and is
memory-bound~\cite{splitwise}, and even large batches leave most of the GPU's
compute unused~\cite{mindgap}. Production systems run the two phases on
separate nodes~\cite{distserve,splitwise,mooncake}. On a decode node, one
lever sets throughput: the batch size. The weights are read once per step
whether the step serves one request or fifty.
% src: splitwise abstract (arXiv 2311.18677): "a compute-intensive prompt
%      computation and a memory-intensive token generation"; "the token
%      generation phase underutilizes compute resources".
%      mindgap abstract (arXiv 2503.08311): "large-batch inference remains
%      memory-bound"; "DRAM bandwidth saturation as the primary bottleneck";
%      "most GPU compute capabilities underutilized". TODO-verify: no
%      utilisation percentage from either PDF is used, on purpose.

The batch size is set by how much KV fits in HBM. A decode step reads every
weight once, so the weights take their full share of HBM before any request
is admitted. For a large model little is left. Qwen2.5-72B in fp8 occupies
76\,GB of a 96\,GB GH200, which leaves about 16\,GB for KV: three requests at
16k tokens. At that batch the compute the GPU can use is bounded below 1\%
(\cref{fig:util}). A longer context shrinks the batch further, in proportion.
% src: paper/model/overhead_model.py Q72 (w = 76.0 GB, b_tok 320 KiB with fp16
%      KV), GH200 96 GB, 4 GB activation reserve: B(16384) = 16e9 /
%      (327680*16384) = 2.98; roofline bound 0.95% (figs/fig-util.tex).
% TODO-verify: growth of context lengths in agentic and reasoning workloads is
%      not claimed until a measured source is fetched.

Host memory is an order of magnitude larger than HBM, and three lines of work
use it to hold KV. The first moves attention off the GPU. CPU attention
offload pays off only when the GPU is weak relative to the CPU: NEO gains up
to $7.5\times$ on a T4 but 14\% on an H100~\cite{neo}. Attention
disaggregation instead runs attention on separate memory-optimised
devices~\cite{lamina}, which is specialised hardware.
% src: neo abstract: "NEO achieves up to 7.5x, 26%, and 14% higher throughput
%      compared to GPU-only approach on T4, A10G, and H100 GPUs".
%      lamina abstract (arXiv 2405.01814 v2): "leverages a collection of cheap,
%      memory-optimized devices for the attention operator".

The second line swaps KV through a staging buffer in HBM.
FlexGen~\cite{flexgen} targets offline, latency-insensitive batches. Pie
places a fixed share of every request's layers in host memory~\cite{pie}, so
every request pays the link every step and none is served purely from HBM.
When the GPU holds enough KV, Pie matches vLLM but does not exceed
it~\cite{pie}. We read this as a bandwidth limit: every offloaded byte crosses
a link far slower than HBM, so offloading extends capacity without adding
bandwidth.
% src: flexgen abstract: "Motivated by the emerging demand for
%      latency-insensitive tasks with batched processing".
%      pie s3.1: "We relocate m layers of the KV cache to the CPU, leaving n-m
%      layers on the GPU. Pie uses per-layer KV cache as the basic management
%      unit"; s5.3 (Fig. 7): "Pie outperforms vLLM at smaller GPU KV cache
%      sizes and achieve similar performance at the largest memory capacity
%      (12GB)". The bandwidth reading is OURS, stated as such.

The third line reads host memory directly from the attention kernel.
DirectKV keeps every request's KV in CPU memory and streams it into the
kernel's shared memory~\cite{directkv}. When HBM suffices, it is slower than
SGLang serving from HBM, and on PCIe it is a capacity extension bounded by
the link~\cite{directkv}.
% src: directkv s6 (l.1030-1033): "During prefill, the manager stores the
%      initial K,V entries from each layer in CPU memory. During decode, the GPU
%      reuses these cached entries directly"; s7.3.1: "SGLang provides the best
%      performance when sufficient HBM is available"; s7.4.3: "on PCIe,
%      DirectKV should be viewed primarily as a capacity-extension mechanism
%      ... remains bounded by CPU-GPU interconnect bandwidth".

Each of these systems sends a given byte of KV through one path, HBM or the
link, and none decides at runtime which requests should take which path.
That decision matters more as links get faster. On a GH200 the GPU reaches
480\,GB of host memory over NVLink-C2C at 900\,GB/s in total, 450\,GB/s in
each direction~\cite{gh200page}; Pie measured 419\,GB/s from host to
device~\cite{pie}. That is several times a PCIe Gen5 link, and enough to move
gigabytes of KV within one decode step. CXL memory adds a further tier of
capacity behind PCIe-class links.
% src: gh200page: "900 gigabytes per second (GB/s) of coherent interface, the
%      superchip is 7X faster than PCIe Gen5"; 480 GB LPDDR5X and 450 GB/s per
%      direction: pie s5.1 ("480GB of LPDDR5X DRAM"), directkv s2 ("900 GB/s,
%      450 GB/s per direction"); pie s2: "419 GB/s from CPU to GPU".
% TODO-verify: CXL 3.x / 4.0 per-x16 bandwidth not stated until the CXL
%      Consortium specification is cited.

We present \sysname{}, a decode engine that serves two classes of request in
the same forward pass (\cref{fig:tiers}c). \emph{Residents} keep their KV in
HBM and read it at HBM bandwidth. \emph{Riders} keep theirs in host memory,
hold no HBM blocks, and stream it layer by layer into a small rotating window
ahead of use, behind the whole batch's compute. Weights and all compute stay
on the GPU. The link becomes a second source of KV bandwidth, and the extra
batch it buys grows with the link and with the share of HBM the weights take
(\cref{fig:potential}).

Riders are not free, and how many pay off depends on the model, the node and
the step time. To our knowledge, \sysname{} is the first serving system that
decides online, every step, how many and which requests ride, from a capacity
model of a measured step time and the link bytes the step can hide. This paper
makes three contributions.
\begin{itemize}\itemsep1pt
\item \textbf{Heterogeneous KV memory management} (\S\ref{sec:mm}). Each
  request is served from exactly one tier, and both tiers run in one forward
  pass. It is built from two blocks: a rolling layer cache with a deferred
  write-back, and step-invariant addressing.
\item \textbf{A capacity model and online admission} (\S\ref{sec:model},
  \S\ref{sec:admission}). A one-time profile and a measured step time decide
  whether riders pay and how many link bytes they may use; admission fills
  that budget every step.
\item \textbf{An implementation in vLLM and its evaluation}
  (\S\ref{sec:impl}--\S\ref{sec:analysis}). An ablation decomposes the cost of
  a second tier, and A100 measurements are projected to GH200 by the
  validated model.
\end{itemize}
